\section{Scope and dependence}\label{sec:scope}

Theorem~\ref{thm:main} closes the classification for one fixed
polynomial, with no restriction on the ambient field or the marked
pair beyond characteristic three. Its proof is short because two
substantial results have already reduced the question to a local
height separation. The new parameter replaces the fixed-point
test by a two-cycle test; it does not supply a new general
equidistribution or local-height theorem.

The argument does not classify arbitrary equal-weight binomials,
nor does it address infinitely intersecting orbits of two fixed
polynomials. It gives a parameter-infinitude criterion, not a
quantitative bound for the finite exceptional parameter sets in
the other cases. There is no change of the iteration clock or
identification of source arithmetic with target Euler factors,
root numbers or a spectral realization.

\paragraph{Preparation and reproducibility.}
All nonexternal mathematical steps are written out in this note;
no experimental enumeration supports an infinite assertion.
The accompanying research record preserves the exact source
versions, full proof review and manuscript-transcription checks.
Lee--Nam is cited at arXiv version~2, dated 11 October 2025.
The Ghioca--Hsia journal metadata were verified at the publisher;
the precise height statement was inspected in Lee--Nam rather
than the original publisher proof. This manuscript was prepared
with AI assistance. Its recorded reviews are internal AI-assisted
checks, not human peer review or global-priority certification.
